\section{Remainders after a cutoff}
\label{sec:remainders}

\subsection{Exponent cutoff and the first possible omitted weight}

Let the target cutoff be an exponent $q$ of the local variable of the result, exclusive. In terms of weights, with $r$ the observable power, a block $z^{r+\sigma}C(L)$ has exponent $(r+\sigma)/|p|$ in the local variable of $y$ (Section~\ref{sec:endpoints}), so the retained set is
\begin{equation}
 I_H=\{\bk\in\N^m:\bk\cdot\bdelta<H\},\qquad H=|p|\,q-r ,
 \label{eq:IH}
\end{equation}
a finite downward-closed set (Lemma~\ref{lem:finite}). This set depends on
the admissible weights, before coefficient cancellation: every index at a
retained weight is included even if its coefficient vanishes. Set
$C^{(r)}_0=1$. When $H>0$ its \emph{one-step boundary} is
\begin{equation}
 \partial I_H=\{\bk+e_i:\bk\in I_H,\ 1\le i\le m\}\setminus I_H ,
 \label{eq:boundary}
\end{equation}
also finite. When $H\le0$, $I_H$ is empty and we define its boundary to be $\{0\}$, since the leading monomial is already omitted. With at least one correction gap, put
\begin{equation}
 \sigma_*=\min\{\bnu\cdot\bdelta:\bnu\in\partial I_H\},\qquad
 C^{(r)}_{\sigma_*}=\sum_{\bk\cdot\bdelta=\sigma_*}C^{(r)}_{\bk},\qquad k_*=\deg C^{(r)}_{\sigma_*}\ (k_*=0\text{ if }C^{(r)}_{\sigma_*}=0).
 \label{eq:frontier-data}
\end{equation}

\begin{lemma}
\label{lem:frontier}
$\sigma_*=\min\{\sigma\in S:\sigma\ge H\}$, and every multi-index
of weight $\sigma_*$ lies in $\partial I_H$. Thus the complete coefficient
at the first \emph{possible} omitted weight is obtained from the boundary
alone, whether or not that coefficient is zero.
\end{lemma}
\begin{proof}
For $H\le0$ the assertion is immediate, with $\sigma_*=0$. For $H>0$, let $\bnu\notin I_H$ have weight $\sigma$. Along a path from $0$ to $\bnu$ that increases one coordinate at a time, the weights increase strictly, and the first point outside $I_H$ is a boundary point of weight $\le\sigma$; hence $\sigma_*\le\sigma$, so $\sigma_*$ is the least weight outside $[0,H)$. If $\sigma=\sigma_*$ the first point outside $I_H$ has weight exactly $\sigma_*$ by minimality, and no later point on the path can have the same weight, so $\bnu$ is that boundary point.
\end{proof}

\begin{theorem}[Remainder after an exponent cutoff]
\label{thm:remainder}
Let $G^{(r)}_H=z^r\sum_{\bk\in I_H}z^{\bk\cdot\bdelta}C^{(r)}_{\bk}(L)$. Then, as $z\to0^+$,
\begin{equation}
 g(y)^r-G^{(r)}_H(y)=z^{r+\sigma_*}C^{(r)}_{\sigma_*}(L)+\oo\bigl(z^{r+\sigma_*}\bigr)
 =\OO\bigl(z^{r+\sigma_*}\M(z)^{k_*}\bigr).
 \label{eq:remainder}
\end{equation}
If $C^{(r)}_{\sigma_*}\ne0$ the bound is attained by its top logarithm.
If $C^{(r)}_{\sigma_*}=0$, the remainder is
$\oo(z^{r+\sigma_*})$. Its first nonzero block, if one exists, occurs at
a later weight and is not identified by this first boundary calculation.
\end{theorem}
\begin{proof}
If $H\le0$, the retained sum is empty and $\sigma_*=0$.
The branch equivalence $g\sim z$ gives
$g^r=z^r(1+\oo(1))$, proving the claim with $C^{(r)}_0=1$.
For $H>0$, Theorem~\ref{thm:convergence} gives $g^r\in\mathcal E$
with the bound \eqref{eq:coefficient-bound}. The omitted multi-indices
have weights at least $\sigma_*$ by Lemma~\ref{lem:frontier}, so the
refined form of Lemma~\ref{lem:tail} applies. When the complete boundary
coefficient is zero, its displayed leading contribution disappears, leaving
the stated little-$o$ estimate.
\end{proof}

The distinction is between the semigroup $S$ of possible weights and the
actual support $\{\sigma\in S:C^{(r)}_\sigma\ne0\}$. A vanishing
coefficient can result from cancellation between indices, from a zero
factor in the coefficient formula, or from $r=0$, when $g^r=1$ and there
is no subsequent nonzero block. If all remaining coefficients are proved
zero, convergence establishes an exact finite expression. A finite run of
zero coefficients alone supplies no such proof.

\begin{corollary}[In the variable of the result]
\label{cor:remainder-y}
With $w$ the local variable of the target ($w=|y-y_0|$ for $p>0$,
and $w=1/|y|$ for $p<0$),
\[
 g(y)^r-G^{(r)}_H(y)=\OO\bigl(w^{(r+\sigma_*)/|p|}\,\M(w)^{k_*}\bigr),\qquad\frac{r+\sigma_*}{|p|}\ge q ,
\]
since $z=(w/|a|)^{1/p}$ for $p>0$, $z=(|a|w)^{1/|p|}$ for $p<0$, and $\M(z)\asymp\M(w)$.
\end{corollary}

For a pure monomial ($m=0$), the inverse observable is exactly $z^r$: the tail is zero if this term is retained and equals $z^r$ otherwise. With correction gaps present, the pair $\bigl((r+\sigma_*)/|p|,k_*\bigr)$ specifies a valid remainder class in the target coordinate, while $z^{r+\sigma_*}C^{(r)}_{\sigma_*}$ is the complete contribution at its first possible omitted weight. If that coefficient vanishes, a sharper order requires computing further weights or proving that the remaining series terminates.

\begin{example}[Resonances]
For $f=x+x^2+x^3$ the gaps are $1$ and $2$; at weight $2$ the indices $(2,0)$ and $(0,1)$ contribute $2$ and $-1$, at weight $3$ the indices $(3,0)$ and $(1,1)$ contribute $-5$ and $+5$, so
\[
 g(y)=y-y^2+y^3+0\cdot y^4-4y^5+14y^6-30y^7+33y^8+\cdots ,
\]
At target cutoff $q=4$ one has $H=3$ and $\sigma_*=3$.
The coarse estimate is $\OO(y^4)$; summing the complete weight-$3$
layer shows that the error is $\oo(y^4)$. To obtain its next order,
include that zero layer in the retained index set and form the new boundary.
At weight $4$, the indices $(4,0)$, $(2,1)$ and $(0,2)$ contribute
$14$, $-21$ and $3$, respectively. Thus
\[
 g(y)-(y-y^2+y^3)=-4y^5+\OO(y^6).
\]
Discarding a zero layer without generating its successors would miss
$(4,0)$ in this calculation. For $f=x+x^2+2x^3$ the block at $y^3$
cancels instead. The cancellation concerns the sum of all contributions
at an exactly equal weight, not each multi-index separately.
\end{example}

\begin{example}[An infinite target and a zero coefficient without resonance]
Let $f(u)=u^{-2}(1+u)$ with $u\to0^+$, so $a=1$, $p=-2$ and
the sole correction gap is $\delta=1$. The positive local inverse and
target coordinate are
\[
 g(y)=\frac{1+\sqrt{1+4y}}{2y}
      =\frac w2+w^{1/2}\sqrt{1+\frac w4},\qquad w=\frac1y\to0^+.
\]
The convergent binomial expansion gives
\[
 g(y)=w^{1/2}+\frac w2+\frac{w^{3/2}}8
                   -\frac{w^{5/2}}{128}+\OO(w^{7/2}).
\]
For the observable $r=1$ and target cutoff $q=2$, one has
$H=|p|q-r=3$. The first possible omitted weight $\sigma_*=3$
would contribute at target exponent $2$, but its coefficient is zero.
There is only one index of each weight here; the zero comes from a factor
in \eqref{eq:master}, not from resonance. The first nonzero omitted term
has weight $4$ and target exponent $5/2$, with coefficient $-1/128$.
This example also shows why a negative leading source power leaves the
correction gaps positive and converts them to target gaps by division by
$|p|$.
\end{example}

\begin{example}[The log example]
Here $m=1$, $\delta=1$, $B=1+L$, so $I_H=\{0,\dots,N\}$ for $H=N+1$, the boundary is $\{N+1\}$, and by \eqref{eq:log-example} the frontier block $y^{N+2}P_{N+1}(\log y)$ has degree exactly $N+1$ with top coefficient $(-1)^{N+1}\Cat_{N+1}$. Hence
\begin{equation}
\begin{aligned}
 g_1(y)&=y+\sum_{n=1}^Ny^{n+1}P_n(\log y)+\OO\bigl(y^{N+2}\M(y)^{N+1}\bigr),\\
 \frac{g_1(y)-y-\sum_{n\le N}y^{n+1}P_n(\log y)}{y^{N+2}(\log y)^{N+1}}&\to(-1)^{N+1}\Cat_{N+1},
\end{aligned}
 \label{eq:log-remainder}
\end{equation}
and in particular after $y-y^2(1+\log y)$ the error is $y^3(2\log^2y+5\log y+3)+\cdots$, not $\OO(y^3)$. The ratio in \eqref{eq:log-remainder} converges slowly (like $1+c/\log y$): for $N=1$ it is $1.89$ at $y=10^{-20}$ and $1.98$ at $y=10^{-100}$.
\end{example}

\subsection{Depth truncation}

When the exponents are symbolic (Section~\ref{sec:boundaries}) their order may be undecidable, but the total degree $|\bk|$ of a multi-index is always known.

\begin{theorem}[Remainder after a depth cutoff]
\label{thm:depth}
Let $N\in\N$ and $m\ge1$, and put
$G^{(r)}_{\le N}=z^r\sum_{|\bk|\le N}z^{\bk\cdot\bdelta}C^{(r)}_{\bk}(L)$,
$\delta_*=\min_i\delta_i$ and $d_*=\max_i\deg B_i$. Then
\begin{equation}
 g(y)^r-G^{(r)}_{\le N}(y)=\OO\bigl(z^{r+(N+1)\delta_*}\M(z)^{(N+1)d_*}\bigr).
 \label{eq:depth-remainder}
\end{equation}
\end{theorem}
\begin{proof}
The omitted terms have $|\bk|\ge N+1$; use \eqref{eq:coefficient-bound} and the count $m^n$ of multi-indices of degree $n$ exactly as in the proof of Lemma~\ref{lem:tail}.
\end{proof}

Depth truncation determines every complete coefficient strictly below
exponent $r+(N+1)\delta_*$. At or beyond that boundary it may retain
terms while omitting larger ones, and a retained coefficient can still
receive contributions from higher depths. For $f=x+x^2+x^4$, depth $1$
gives $y-y^2-y^4$, whereas the inverse begins
$y-y^2+2y^3-6y^4+\OO(y^5)$. Thus the displayed $y^4$ term is both
smaller than the omitted $2y^3$ term and only part of the true coefficient
at $y^4$. The depth remainder $\OO(y^3)$ is nevertheless correct.
Depth is an inclusive bound on $|\bk|$; it is neither an exponent cutoff
nor a count of resolved nonzero blocks. For $m=0$ there are no correction
indices and the inverse observable is the exact monomial $z^r$.
